\documentclass[aspectratio=169]{beamer}
\usetheme{Madrid}
\usecolortheme{beaver}

\usepackage[UTF8,fontset=fandol]{ctex}
\usepackage{amsmath, amssymb}
\usepackage{booktabs}
\usepackage{ragged2e}
\usepackage{array}
\usepackage{graphicx}
\usepackage{tikz}
\usetikzlibrary{positioning, arrows.meta, shapes.geometric, calc, fit}

\newcommand{\framebottomnote}[1]{%
    \vfill
    \par\noindent
    {\tiny\color{black!65}\parbox{0.96\linewidth}{#1}\par}%
    \vspace*{-0.75cm}%
}

\newcolumntype{L}[1]{>{\RaggedRight\arraybackslash}p{#1}}

\title[Helplessness 机制的初步形式化]{Digital Friction 中 Helplessness 机制的初步形式化}
\date{2026-04-20}

\begin{document}

\begin{frame}
    \titlepage
\end{frame}

\section{问题定义}

\begin{frame}[t]{相关论文如何理解 Scope}
\small
\begin{columns}[T]
\begin{column}{0.48\textwidth}
\textbf{Learned Helplessness in Humans: Critique and Reformulation}

\textit{Lyn Y. Abramson, Martin E. P. Seligman, John D. Teasdale}
\begin{itemize}
    \item Abramson, Seligman, and Teasdale (1978) 认为，人类的 helplessness 不能只看其结果是否不可控，更关键的是\textbf{人怎样解释这次 noncontingency}
    \item \mbox{\textbf{scope决定generality}}，如果个体把失败理解成“这一类事我都不行”，无助就更容易从当前任务扩散出去
\end{itemize}
\end{column}

\begin{column}{0.48\textwidth}
\textbf{Situation Similarity and the Generalization of Learned Helplessness}

\textit{Marika Tiggemann, A. H. Winefield}
\begin{itemize}
    \item Tiggemann and Winefield (1978) 发现，learned helplessness 在相似任务上明显出现，但在不相似的任务上几乎不泛化
    \item 这说明 helplessness 的跨任务扩散不是无边界的，而是受\textbf{场景相似度}约束
\end{itemize}
\end{column}
\end{columns}

\vspace{0.15cm}
\centering
\small
\[
\text{Failure} \rightarrow \text{Attribution Scope} \rightarrow \text{Task-Family Spillover}
\]
\end{frame}

\section{形式化表达}

\begin{frame}[c]{Gaussian为什么被选中作为scope的形式化更新方式}
\small
\begin{center}
\begin{minipage}[t][0.78\textheight][t]{0.84\textwidth}
\justifying
\begin{itemize}
    \item 如果一次数字事件失败的 Scope 归因也会扩散，那它应该优先影响相似任务，而不是平均地扩散到所有任务。由于目前缺少针对 Scope 的成熟形式化研究，这里\textbf{借鉴了人类泛化机制中的距离衰减写法。}
    \item Lucas et al. (2015) 和 Wu et al. (2018) 在人类泛化建模里，正是用高斯过程来表达这种“按距离衰减”的关系
    \item Lee et al. (2021) 进一步指出，泛化梯度常接近高斯形状；按心理距离缩放后，又和 Shepard 的指数衰减规律一致，所以 Gaussian 在这里既符合直觉，也有认知建模上的理论支撑
\end{itemize}

\vspace{0.1cm}
\textbf{区分：}我们借用的是距离衰减形式，不是把 scope 和人类泛化当成同一个变量。

\vfill
{\tiny \color{black!60}\textbf{人类泛化机制：}人会把一次经验推广到相似情境中，而且这种推广会随着相似度下降而减弱。}
\end{minipage}
\end{center}
\end{frame}

\begin{frame}[t]{Gaussian Kernel 版本}
\scriptsize
\justifying
\[
A_t=\beta\, U_t\, a_t
\]

\[
d(i,j)=1-S_{ij}
\]

\[
w_{ij}=\exp\!\left(-\frac{d(i,j)^2}{2\sigma(g_t)^2}\right), \qquad j\neq i
\]

\[
\hat w_{ij}=\frac{w_{ij}}{\sum_{k\neq i} w_{ik}}
\]

\[
\Delta SE^{\text{spill}}_{j,t}
=
-A_t\, \hat w_{ij}
\]

\vspace{0.1cm}
\raggedright
\textbf{标注：}当前代码里，\(a_t\) 是 LLM 输出的 \texttt{scope amplitude}；\(\sigma\) 先固定；\(\hat w_{ij}\) 是归一化后的 Gaussian 权重。
\end{frame}

\section{实验接入}

\begin{frame}[t]{在当前实验中如何加入}
\scriptsize
\begin{columns}[T]
\begin{column}{0.54\textwidth}
\textbf{当前实验里的流程}

\vspace{0.05cm}
\[
\text{数字任务失败事件}
\]
\[
\downarrow
\]
\[
\text{LLM 输出 scope amplitude } a_t
\]
\[
\downarrow
\]
\[
\text{若 attribution 有效且过阈值，计算 } A_t=\beta U_t a_t
\]
\[
\downarrow
\]
\[
\text{结合相似度矩阵 } S_{ij} \text{ 与 Gaussian 权重分配}
\]
\[
\downarrow
\]
\[
\text{扩散到其他 task family 的 self-efficacy}
\]
\[
\downarrow
\]
\[
\text{进一步影响后续 decision / attempt / outcome}
\]
\end{column}

\begin{column}{0.40\textwidth}
\textbf{当前实现的版本}

\[
\textbf{单参数 Gaussian spillover}
\]

\begin{itemize}
    \setlength{\itemsep}{0.1em}
    \item LLM 只输出 scope amplitude
    \item \(\sigma\) 先固定，不让 LLM 参与；后续再决定是否加入可变 \(\sigma\)
    \item Gaussian 权重会做归一化
\end{itemize}

\vspace{0.02cm}
\textbf{当前代码使用的 fixed similarity matrix}

\centering
\renewcommand{\arraystretch}{0.92}
\begin{tabular}{@{}lcccc@{}}
\toprule
 & L & A & P & H \\
\midrule
L & 1.0 & 0.3 & 0.8 & 0.2 \\
A & 0.3 & 1.0 & 0.2 & 0.7 \\
P & 0.8 & 0.2 & 1.0 & 0.2 \\
H & 0.2 & 0.7 & 0.2 & 1.0 \\
\bottomrule
\end{tabular}

\vspace{0.03cm}
\raggedright
{\tiny
L = 登录验证, A = 挂号预约\\
P = 支付结算, H = 健康信息查询
}

\vspace{0.02cm}
{\tiny
直觉上，L--P 更近，A--H 更近，跨组传播更弱。
}
\end{column}
\end{columns}
\end{frame}

\section{主线补充}

\begin{frame}[t]{A Bayesian formulation of behavioral control}
\small
\begin{itemize}
    \item \textbf{第一层：outcome entropy}。看一个动作会不会稳定导向少数结果，如果越稳定，控制就越高（更偏向于单个动作的稳定性）
    \item \textbf{第二层：controllably achievable outcomes}。看不同动作是不是能稳定带来不同结果，而不是多个动作的情况下，最后结果差不多（但这一层还不涉及这些结果是否真正有价值）
    \item \textbf{第三层：controllably achievable reinforcement}。最关键的是那些对个体有价值的结果里，有多少是可以通过行为可靠控制的
\end{itemize}

\vspace{0.08cm}
\textbf{论文最后更看重的是第三层，因为它更接近 helplessness 的心理实质：}

\[
\text{我在乎的结果，到底是不是受我行为控制？}
\]
\end{frame}

\begin{frame}[c]{A Bayesian formulation of behavioral control}
\small
\centering
\begin{minipage}{0.86\textwidth}
\begin{itemize}
    \item 论文形式化的是个体对环境可控性的先验信念，即 \textbf{controllability prior}
    \item 一旦这种先验变得悲观，它就会进入后续任务，影响 action-outcome 预测、动作价值差异和探索意愿
    \item 被动、少尝试和更容易放弃，更接近于低 controllability prior 在新任务中的行为后果
\end{itemize}
\end{minipage}

\vspace{0.18cm}
\centering
\[
\text{过去经验}
\rightarrow
\text{可控性先验}
\rightarrow
\text{当前任务中的结果预测与动作价值判断}
\rightarrow
\text{更被动的行为}
\]
\end{frame}

\begin{frame}[c]{核心预测分布}
\scriptsize
\centering
\[
P(r\mid N_t,a,\alpha)
=
\frac{\alpha}{\alpha+N_t(a)}B(r)
+
\frac{1}{\alpha+N_t(a)}N_t(r,a)
\]

\vspace{0.1cm}
\begin{minipage}{0.84\textwidth}
\begin{itemize}
    \item 这个式子表示：\textbf{如果现在再做一次动作 \(a\)，结果 \(r\) 的预测概率是多少}
    \item 预测不是只由眼前经验决定，也不是只由原有先验决定，而是由两者共同决定
    \item \(B(r)\)：还没什么经验时，对结果 \(r\) 的默认预期
    \item \(N_t(r,a)\) 和 \(N_t(a)\)：到目前为止，动作 \(a\) 一共试了多少次，其中结果 \(r\) 出现了多少次
    \item \(\alpha\)：先验权重；\(\alpha\) 越大，少量新经验越难立刻推翻原来的判断
\end{itemize}
\end{minipage}
\end{frame}

\end{document}
